\documentclass[12 pt]{article}

\usepackage{stoversymb}
\usepackage[margin=0.75in, top=0.875in, bottom=0.875in]{geometry}
\usepackage{amsmath,amsfonts,amssymb,url,multicol,graphicx,tikz,soul,changepage}
%===makes urls render well===
\usepackage{lmodern}
\usepackage[T1]{fontenc}
%============================
\usepackage{wasysym} % smileys
\usepackage[inline]{enumitem}
\usepackage[many]{tcolorbox}
\tcbuselibrary{breakable}

%\everymath{\displaystyle}

\makeatletter
\renewcommand*\env@matrix[1][\arraystretch]{%
	\edef\arraystretch{#1}%
	\hskip -\arraycolsep
	\let\@ifnextchar\new@ifnextchar
	\array{*\c@MaxMatrixCols c}}
\makeatother

%\setenumerate{itemsep=0.25in}
%\setlist[enumerate,1]{label=\arabic*., itemsep=0.5in}
%\setlist[enumerate,2]{label={(\alph*)}, itemsep=0.8in}
%\setlist[enumerate,2]{leftmargin=0.5in}
%\setlist[enumerate,3]{label={\roman*.}}
\setlist[itemize,1]{label=$\circ$,topsep=0mm,itemsep=0mm}

\newcommand{\truefalse}[1]{#1\hfill\rule[-1mm]{220pt}{0.75pt}}
\newcommand{\hint}[1]{\vspace{3mm}\textbf{Hint}: #1}
\newcommand{\note}[1]{\textbf{Note}: #1}
\newcommand{\infsum}[3]{\sum_{{#1}={#2}}^\infty {#3}}
\newcommand{\axes}[1]{\begin{center}\includegraphics[scale=#1]{3DAxes}\end{center}}
\newcommand{\comps}[1]{\langle #1_1,#1_2,#1_3\rangle}
\newcommand{\compslong}[3]{\langle #1, #2, #3\rangle}
\newcommand{\ijk}[2]{#1\vect{#2}}
%\newcommand{\pmat}[1]{\begin{pmatrix}[1.25]#1\end{pmatrix}}
%\newcommand{\bmat}[1]{\begin{bmatrix}[1.25]#1\end{bmatrix}}
\newcounter{Exx}
\newcommand{\exbox}[1]{\stepcounter{Exx}\begin{tcolorbox}[breakable,arc=2pt,boxrule=1pt,right=6mm,top=3mm,bottom=6mm]\noindent\textit{\large\bfseries Example \theExx:\\[3mm]}~#1\end{tcolorbox}}
\newcommand{\exit}[1]{\vspace{4.5mm}\noindent\textit{Example #1:}}
\newcommand{\exbf}[1]{\vspace{1.5mm}\noindent\textbf{Example #1:}}
\newcommand{\LHS}{\text{LHS}}
\newcommand{\RHS}{\text{RHS}}
\newcommand{\ans}{\noindent\underline{Ans:}~}

\newenvironment{mypmatrix}[1]{\renewcommand{\arraystretch}{#1}\begin{pmatrix}}{\end{pmatrix}}
\newenvironment{mybmatrix}[1]{\renewcommand{\arraystretch}{#1}\begin{bmatrix}}{\end{bmatrix}}
\newcommand{\justify}{\ul{Justify your claim.}}

\newcommand{\pmat}[1]{\begin{mypmatrix}{1.25}#1\end{mypmatrix}}
\newcommand{\bmat}[1]{\begin{mybmatrix}{1.25}#1\end{mybmatrix}}

\makeatletter
\renewcommand*\env@matrix[1][*\c@MaxMatrixCols c]{%
	\hskip -\arraycolsep
	\let\@ifnextchar\new@ifnextchar
	\array{#1}}
\makeatother

\newcommand{\pmatgrid}[2]{\renewcommand{\arraystretch}{1.25}\begin{pmatrix}[#1] #2\end{pmatrix}}
\newcommand{\bmatgrid}[2]{\renewcommand{\arraystretch}{1.25}\begin{bmatrix}[#1] #2\end{bmatrix}}
\newcommand{\warnbox}[1]{\begin{adjustwidth}{0.375in}{0.375in}\begin{tcolorbox}[arc=0pt, colback=red!5!white, colframe=red, boxrule=0.5pt, parbox=false]{\parskip=3mm \parindent=0mm \red{\textbf{\large\itshape Warning!}\par #1}}\end{tcolorbox}\end{adjustwidth}}
\newcommand{\greentheorembox}[2]{\begin{adjustwidth}{0.375in}{0.375in}\begin{tcolorbox}[arc=0pt, colback=darkgreen!10!white, colframe=darkgreen, boxrule=0.5pt, parbox=false]{\parskip=3mm \parindent=0mm \green{\textbf{\large\itshape #1:}\par #2}}\end{tcolorbox}\end{adjustwidth}}
\newcounter{FactBox}
\newcommand{\factbox}[1]{\stepcounter{FactBox}\begin{adjustwidth}{0.375in}{0.375in}\begin{tcolorbox}[arc=0pt, colback=darkgreen!10!white, colframe=darkgreen, boxrule=0.5pt, parbox=false]{\parskip=3mm \parindent=0mm \green{\textbf{\large\itshape Fact \theFactBox:}\par #1}}\end{tcolorbox}\end{adjustwidth}}

\newcommand{\arrow}[1]{\xrightarrow{\hspace*{#1}}}

\setlength\parskip{0.75mm}	

\newcounter{ThmBox}
\newcommand{\thmbox}[1]{\stepcounter{ThmBox}\vspace{6mm}\noindent\fbox{\parbox{\linewidth}{~~(\theThmBox)~\hfill #1 \hfill\hfill}}}
\newcommand{\why}[1]{\begin{adjustwidth}{0.5in}{0.5in}{\par\vspace{3mm}\noindent\green{\textbf{Why should this be true? }#1}}\qed\end{adjustwidth}\vspace{6mm}}

\usepackage{tikz}
\usetikzlibrary{decorations.pathreplacing}

\newcommand{\tikzmark}[1]{\tikz[overlay,remember picture] \node[baseline] (#1) {};}

\tikzset{My Node Style/.style={midway, right, xshift=1mm, align=left, font=\small, draw=none, thin, text=black}}

\newcommand\VerticalBrace[4][]{%
	% #1 = draw options
	% #2 = top mark
	% #2 = bottom mark
	% #4 = label
	\begin{tikzpicture}[overlay,remember picture]
	\draw[decorate,decoration={brace, amplitude=0.5ex}, #1] 
	([yshift=1ex,xshift=3in]#2.north east)  -- ([yshift=0ex,xshift=3in]#3.south east)
	node[My Node Style] {#4};
	\end{tikzpicture}
}

\newcommand\LeftVertBrace[4][]{%
	% #1 = draw options
	% #2 = top mark
	% #2 = bottom mark
	% #4 = label
	\begin{tikzpicture}[overlay,remember picture]
	\draw[decorate,decoration={brace, amplitude=0.5ex}, #1] 
	([yshift=-3ex,xshift=-1.71in]#2.north west)  -- ([yshift=4ex,xshift=1.625in]#3.south west)
	node[My Node Style,below,xshift=-0.5in,yshift=0.5in] {#4};
	\end{tikzpicture}
}

\newcommand{\matA}{\pmat{1 & 2 & 3 & 4 & 5 & 6 \\ 0 & -1 & 1 & -3 & 4 & -5 \\ 1 & 0 & 0 & 1 & -2 & -4 \\ 0 & 0 & 1 & 1 & 2 & -7}}
\newcommand{\matAred}{\pmat{1 & 0 & 0 & 0 & -\dfrac{11}{8} & \dfrac{3}{8} \\[4.5mm] 0 & 1 & 0 & 0 & \dfrac{1}{2} & \dfrac{31}{2} \\[4.5mm] 0 & 0 & 1 & 0 & \dfrac{21}{8} & -\dfrac{21}{8} \\[4.5mm] 0 & 0 & 0 & 1 & -\dfrac{5}{8} & -\dfrac{35}{8}}}

\graphicspath{ {./../img/} }
\DeclareGraphicsExtensions{.pdf}

\begin{document}
\begin{center}
\Large{\textbf{Vector Spaces}}
\end{center}
Recently, I mentioned that I'd begin using the phrase ``vector space'' but wouldn't formally define it in class. To alleviate the burden caused, I've decided to give you a handout on the topic instead.

Here's the definition

\vspace{3mm}

\begin{adjustwidth}{6mm}{6mm}
	\noindent\textbf{Definition 1:} A \textit{vector space} is a nonempty set $V$ of objects (called \textit{vectors}) and a pair of operations---called \textit{addition} (denoted ``$+$'') and \textit{scalar multiplication} (denoted by juxtaposition, i.e. $c\vect{u}$ for the scalar multiple of $\vect{u}$ by $c$)---which together will always satisfy the following ten axioms (for all scalars $c$ and $d$ and for all vectors $\vect{u}$, $\vect{v}$, and $\vect{w}$ in $V$):
	\begin{enumerate}[label=(\roman*),leftmargin=16mm,itemsep=4.5mm,rightmargin=3in,topsep=6mm]
		\item $\vect{u}+\vect{v}\in V$\hspace{-2.75in}\tikzmark{top 1}\tikzmark{bottom 1}\tikzmark{vbottom 1}
		\item $c\vect{u}\in V$\hspace{-2.53in}\tikzmark{top 2}\tikzmark{bottom 2}
		\item \tikzmark{top 3}there exists a zero vector $\vect{0}\in V$ which satisfies $\vect{u}+\vect{0}=\vect{u}$\tikzmark{vtop 1}
		\item $\vect{u}+\vect{v}=\vect{v}+\vect{u}$
		\item $\vect{u}+\left(\vect{v}+\vect{w}\right)=\left(\vect{u}+\vect{v}\right)+\vect{w}$
		\item for each $\vect{u}\in V$, there exists a vector $-\vect{u}\in V$ such that $\vect{u}+\left(-\vect{u}\right)=\vect{0}$
		\item $c\left(\vect{u}+\vect{v}\right)=c\vect{u}+c\vect{v}$
		\item $(c+d)\vect{u}=c\vect{u}+d\vect{u}$
		\item $c\left(d\vect{u}\right)=\left(c d\right)\vect{u}$
		\item $1\vect{u}=\vect{u}$\hspace{-0.525in}\tikzmark{bottom 3}
	\end{enumerate}

	\VerticalBrace[ultra thick, red]{top 1}{bottom 1}{\red{We can add things!}}
	\VerticalBrace[ultra thick, red]{top 2}{bottom 2}{\red{We can scalar multiply}}
	\VerticalBrace[ultra thick, red]{top 3}{bottom 3}{\parbox{3in}{\red{Addition and scalar multiplication work the way we want them to!}}}
	\LeftVertBrace[ultra thick, red]{vtop 1}{vbottom 1}{\parbox{0.5in}{\scriptsize{\red{These should look like the subspace axioms!}}}}
\end{adjustwidth}

\vspace{-6mm}

Over the next week or so, my plan is to revisit this handout $\geq 1$ time(s) to give examples, details for verifying examples, etc., of vector spaces. Until then, just remember the mantra:

\begin{center}
	\huge{When you see ``\red{...$n$-dimensional vector space $V$...}'' just replace it with ``\red{...subspace of $\Reals^n$...}''!}
\end{center}

%\subsection*{\hspace{-6.5mm}Column Space}
%\noindent\ul{What is it?} 
%
%The column space of $\sansA$ (denoted $\col(\sansA)$) is the span of the columns of $\sansA$:
%\begin{equation}
%	\label{colSpaceEq}
%	\col(\sansA)=\vsspan\{\text{the columns of $\sansA$}\}.
%\end{equation}
\end{document}